%% lampiran/A-detail-fitur.tex
%% Lampiran A — Detail Fitur Time-Domain dan Frequency-Domain

\chapter{Detail Fitur Time-Domain dan Frequency-Domain}
\label{lamp:fitur}

Lampiran ini merinci definisi numerik vektor \emph{Health Indicator} (HI)
yang memasuki model pada Bab~\ref{bab:metodologi-umum}--\ref{bab:prognostik}.
Implementasi mengikuti modul \path{mxlstm/data/hi.py} pada repositori
publik (Lampiran~\ref{lamp:kode}).

Setiap rekaman getaran memuat dua kanal akselerometer (horizontal lalu
vertikal, sesuai urutan kolom pada berkas CSV adaptor dataset).
Untuk setiap kanal $c \in \{0,1\}$ diekstraksi sembilan skalar
\emph{time-domain} dan sembilan skalar turunan spektral berbasis PSD
Welch, sehingga
\begin{equation}\label{eq:lampA_dim}
  F = 2 \times (9 + 9) = 36.
\end{equation}
\noindent Konvensi penamaan fitur dalam kode dan pada tabel SHAP/IG
Bab~\ref{bab:prognostik} mengikuti pola
\texttt{td\_c\textit{c}\_}\textit{nama} dan
\texttt{fd\_c\textit{c}\_}\textit{nama}, dengan indeks kanal
$c=0$ untuk horizontal dan $c=1$ untuk vertikal.

Ringkasan energi pada lima pita frekuensi seragam hingga Nyquist
(sebagaimana diuraikan pada Bab~\ref{bab:metodologi-umum})
menggambarkan prinsip dekomposisi spektral; secara numerik, energi pita
pada Lampiran ini dihitung sebagai jumlah kontribusi bin PSD Welch
pada rentang frekuensi yang sama lebar, sehingga konsisten dengan
implementasi pelatihan.

\section{Notasi Bersama}
\label{lamp:fitur_notasi}

Misalkan $(x_1, x_2, \ldots, x_N)$ menyatakan sampel sinyal satu kanal
pada satu rekaman, dengan $N$ panjang rekaman dan
$f_s = 25{,}6$~kHz frekuensi \emph{sampling}.
Diberikan $\mu = \frac{1}{N}\sum_{i=1}^{N} x_i$,
$\sigma^2 = \frac{1}{N}\sum_{i=1}^{N}(x_i - \mu)^2$,
dan $\sigma = \sqrt{\sigma^2 + \varepsilon}$ dengan
$\varepsilon > 0$ kecil untuk stabilitas numerik.
Notasi $\mathbb{E}[\cdot]$ pada skewness dan kurtosis di bawah ini
merujuk pada momen empiris terhadap distribusi sampel yang sama.

\section{Fitur Time-Domain}
\label{lamp:fitur_time}

Tabel~\ref{tab:lampA_time} merinci sembilan fitur \emph{time-domain}
per kanal. Urutan vektor pada keluaran program mengikuti baris tabel.

\begin{table}[ht]
  \centering
  \caption{Definisi sembilan fitur \emph{time-domain} per kanal.}
  \label{tab:lampA_time}
  \footnotesize
  \begin{tabularx}{\textwidth}{@{}c l X@{}}
    \toprule
    \textbf{No.} & \textbf{ID (\texttt{td\_c$c$\_\dots})} & \textbf{Definisi} \\
    \midrule
    1 & \texttt{rms} &
    $\mathrm{RMS} = \sqrt{\frac{1}{N}\sum_{i=1}^{N} x_i^2}$. \\
    2 & \texttt{peak} &
    $\mathrm{Peak} = \max_{1 \le i \le N} |x_i|$. \\
    3 & \texttt{kurtosis} &
    Kurtosis Pearson
    $\mathbb{E}\bigl[(x_i-\mu)^4\bigr]/\sigma^4$
    (parameter \texttt{fisher=False} pada \texttt{scipy.stats.kurtosis},
    sehingga untuk sinyal Gaussian mendekati $3{,}0$). \\
    4 & \texttt{skewness} &
    Skewness
    $\mathbb{E}\bigl[(x_i-\mu)^3\bigr]/\sigma^3$. \\
    5 & \texttt{crest} &
    \emph{Crest factor} $\mathrm{Peak}/(\mathrm{RMS}+\varepsilon)$. \\
    6 & \texttt{shape} &
    \emph{Shape factor}
    $\mathrm{RMS}/\bigl(\frac{1}{N}\sum_{i=1}^{N}|x_i| + \varepsilon\bigr)$. \\
    7 & \texttt{impulse} &
    \emph{Impulse factor}
    $\mathrm{Peak}/\bigl(\frac{1}{N}\sum_{i=1}^{N}|x_i| + \varepsilon\bigr)$. \\
    8 & \texttt{margin} &
    \emph{Margin factor}
    $\mathrm{Peak}/\bigl(\bigl(\frac{1}{N}\sum_{i=1}^{N}\sqrt{|x_i|}\bigr)^2 + \varepsilon\bigr)$. \\
    9 & \texttt{variance} &
    Variansi populasi
    $\sigma^2 = \frac{1}{N}\sum_{i=1}^{N}(x_i-\mu)^2$. \\
    \bottomrule
  \end{tabularx}
\end{table}

\section{Fitur Frequency-Domain (PSD Welch)}
\label{lamp:fitur_freq}

Spektral dihitung dengan metode Welch (\texttt{scipy.signal.welch}):
panjang segmen FFT
$n_{\mathrm{perseg}} = \min(2048,\, N)$,
jendela default bibliotek, dan vektor frekuensi $f_\ell$ serta PSD
diskrit $P_\ell$ pada bin yang dihasilkan.
Didefinisikan bobot ternormalisasi
\begin{equation}\label{eq:lampA_psd_norm}
  \tilde{P}_\ell = \frac{P_\ell}{\sum_{j} P_j + \varepsilon},
  \qquad
  \sum_{\ell} \tilde{P}_\ell \approx 1.
\end{equation}

\noindent Empat statistik spektral global (baris 1--4 pada
Tabel~\ref{tab:lampA_freq}) menggunakan $P_\ell$ mentah untuk
\emph{mean frequency} dan \emph{RMS frequency}, serta $\tilde{P}_\ell$
untuk pusat massa spektral dan entropi spektral.
Lima energi pita (baris 5--9) merupakan jumlah $P_\ell$ pada bin yang
jatuh ke dalam lima sub-interval seragam
$[0,\, f_s/2]$:
\begin{equation}\label{eq:lampA_band}
  \mathcal{B}_k =
  \Bigl[f_{k-1},\, f_k\Bigr),
  \quad
  f_k = k \cdot \frac{f_s/2}{5},
  \quad
  k = 1,\ldots,5,
\end{equation}
\begin{equation}\label{eq:lampA_band_energy}
  E_k = \sum_{f_\ell \in \mathcal{B}_k} P_\ell .
\end{equation}

\begin{table}[ht]
  \centering
  \caption{Definisi sembilan fitur frekuensi per kanal (Welch PSD).}
  \label{tab:lampA_freq}
  \footnotesize
  \begin{tabularx}{\textwidth}{@{}c l X@{}}
    \toprule
    \textbf{No.} & \textbf{ID (\texttt{fd\_c$c$\_\dots})} & \textbf{Definisi} \\
    \midrule
    1 & \texttt{centroid} &
    Frekuensi pusat massa spektrum terbobot $\tilde{P}$:
    $\sum_{\ell} f_\ell\, \tilde{P}_\ell$. \\
    2 & \texttt{entropy} &
    Entropi Shannon
    $-\sum_{\ell : \tilde{P}_\ell > 0}
    \tilde{P}_\ell \ln(\tilde{P}_\ell + \varepsilon)$. \\
    3 & \texttt{mean\_freq} &
    \emph{Mean frequency}
    $\sum_{\ell} f_\ell P_\ell / (\sum_{j} P_j + \varepsilon)$. \\
    4 & \texttt{rms\_freq} &
    \emph{RMS frequency}
    $\sqrt{\sum_{\ell} f_\ell^2 P_\ell / (\sum_{j} P_j + \varepsilon)}$. \\
    5--9 & \texttt{band0}--\texttt{band4} &
    Energi pita $E_k$ pada Persamaan~\eqref{eq:lampA_band_energy}
    untuk $k = 1,\ldots,5$ (berkas kode memakai indeks nol:
    \texttt{band0} $\equiv E_1$, \ldots, \texttt{band4} $\equiv E_5$). \\
    \bottomrule
  \end{tabularx}
\end{table}

\section{Pasca-Ekstraksi pada Rantai Pelatihan}
\label{lamp:fitur_post}

Setelah vektor $36$-dimensi terbentuk per rekaman, diterapkan
normalisasi \emph{Min--Max} per fitur yang di-\emph{fit} hanya pada
bantalan pelatihan, diikuti pemulusan eksponensial sepanjang indeks
waktu dengan koefisien $\alpha_{\mathrm{smooth}} = 0{,}10$
(sebagaimana pada Bab~\ref{bab:metodologi-umum}).
Parameter normalisasi disimpan per fitur agar dapat diterapkan identik
pada validasi dan uji tanpa penyimpangan statistik antar-set.
Urutan operasi dan parameter persis tercatat pada kelas
\texttt{HIPipeline} dalam \path{mxlstm/data/hi.py}.
